%!TEX program = uplatex
\RequirePackage{plautopatch}

\documentclass[uplatex,dvipdfmx]{jlreq}

% 余白調整（1ページ収め）
\usepackage[top=18mm,bottom=20mm,left=20mm,right=20mm]{geometry}

\usepackage{amsmath,amssymb}
\usepackage{url}

\pagestyle{empty}

% 行間少し詰める
\renewcommand{\baselinestretch}{0.96}

%--------------------------------
% 講演マクロ
%--------------------------------

\newcommand{\talk}[3]{%
\noindent
#1\quad #2\par
\hspace*{2em}#3\par\smallskip
}

%--------------------------------
% タイトル情報
%--------------------------------

\newcommand{\EwMnumber}{第21回}
\newcommand{\EwMtitle}{実解析への誘い\\[2mm]\large 実解析的方法を使いこなそう}
\newcommand{\EwMdate}{2001年10月26日（金）14:30～18:00 ～ 10月27日（土）10:30～17:00}
\newcommand{\EwMplace}{東京都文京区春日1-13-27 中央大学理工学部5号館 5533号室}

%--------------------------------
% タイトル表示
%--------------------------------

\newcommand{\EwMtitleblock}{

\vspace*{-5mm}

\begin{center}

{\Large ENCOUNTER with MATHEMATICS}

\vspace{2mm}

{\large 数学との遭遇}

\vspace{4mm}

{\Large \bfseries \EwMnumber\ \EwMtitle}

\vspace{4mm}

\EwMdate

\vspace{2mm}

於：\EwMplace

\end{center}

\vspace{2mm}

}

\begin{document}

\EwMtitleblock

%--------------------------------
% プログラム
%--------------------------------

\section*{プログラム}

\subsection*{10月26日（金）}

\talk
{14:30～15:50}
{実解析的方法とはどのようなものか}
{新井 仁之 氏（東大・数理）}

\talk
{16:30～18:00}
{最大関数と Littlewood–Paley 関数}
{宮地 晶彦 氏（東京女子大・文理）}

\subsection*{10月27日（土）}

\talk
{10:30～12:00}
{ストリッカーズ評価とその応用}
{小澤 徹 氏（北大・理）}

\talk
{14:00～15:30}
{フラクタル上の解析学\\[1mm]自己相似集合上のラプラシアン}
{木上 淳 氏（京大・情報）}

\talk
{15:50～17:00}
{ブラウン運動と実解析\\[1mm]実解析のための確率論入門}
{新井 仁之 氏（東大・数理）}

\bigskip

17:20～　懇親会（ワイン・パーティー）

%--------------------------------
% 案内文
%--------------------------------

\bigskip

別紙の趣旨に沿った集会の第21回を以上のような予定で開催いたします。  
非専門家向けに入門的な講演をお願い致しました。  
多く方々のご参加をお待ちしております。  
講演者による講演内容へのご案内を別紙に添付いたしますのでご覧下さい。

%--------------------------------
% 連絡先
%--------------------------------

\section*{連絡先}

112-8551 東京都文京区春日1-13-27 中央大学理工学部数学教室 

tel : 03-3817-1745  

ENCOUNTER with MATHEMATICS : e-mail : encmath@math.chuo-u.ac.jp  
homepage : \url{http://www.math.chuo-u.ac.jp/ENCwMATH}  

三松 佳彦 : yoshi@math.chuo-u.ac.jp / 高倉 樹 : takakura@math.chuo-u.ac.jp

\end{document}
t